\documentclass[11pt]{article}
\usepackage{amsmath,amssymb,amsthm}
\usepackage{algorithm,algorithmic}
\usepackage{hyperref}
\usepackage{enumitem}
\usepackage{bm}
\usepackage{color}
\newtheorem{theorem}{Theorem}
\newtheorem{lemma}{Lemma}
\newcommand{\prox}{\operatorname{prox}}
\newcommand{\E}{\mathbb{E}}
\newcommand{\norm}[1]{\left\lVert #1\right\rVert}
\newcommand{\inner}[2]{\left\langle #1,#2\right\rangle}
\newcommand{\R}{\mathbb{R}}

\begin{document}

\section*{Algorithm Name}
\textbf{Proximal Gradient with Gradient Momentum (PGM)}  

The method solves the composite convex problem  

\[
\min_{x\in\R^{d}}\; \Phi(x)\quad\text{with}\quad \Phi(x)=f(x)+g(x),
\]

where $f$ is smooth (Lipschitz gradient) and $g$ is closed, proper and convex (possibly nonsmooth).

\vspace{0.2cm}
\section*{Mathematical Formulation}
Given a stepsize $\eta>0$ and a momentum parameter $\beta\in[0,1)$, the iterates are defined as  

\[
\begin{aligned}
\text{(Momentum on the gradient)}\qquad
&\tilde\nabla f_{k}= \nabla f(x_k)+\beta\bigl(\nabla f(x_k)-\nabla f(x_{k-1})\bigr), \tag{1}\\[4pt]
\text{(Proximal step)}\qquad
&x_{k+1}= \prox_{\eta g}\!\bigl(x_k-\eta \tilde\nabla f_{k}\bigr), \qquad k\ge 0,
\end{aligned}
\]
with a dummy initialization $\nabla f(x_{-1})\equiv 0$ (or $x_{-1}=x_0$).  
The proximal operator is  

\[
\prox_{\eta g}(z)=\arg\min_{u\in\R^{d}}\Bigl\{g(u)+\frac{1}{2\eta}\norm{u-z}^{2}\Bigr\}.
\]

\vspace{0.2cm}
\section*{Pseudocode}
\begin{algorithm}[H]
\caption{Proximal Gradient with Gradient Momentum (PGM)}
\begin{algorithmic}[1]
\Require stepsize $\eta>0$, momentum $\beta\in[0,1)$, initial point $x_{0}\in\R^{d}$
\State Set $x_{-1}\gets x_{0}$ \Comment{needed for the first momentum term}
\For{$k=0,1,2,\dots$}
    \State Compute stochastic (or deterministic) gradient $g_{k}\gets\nabla f(x_{k})$
    \State Compute momentum‑augmented gradient 
    \[
    \tilde g_{k}\gets g_{k}+\beta\bigl(g_{k}-\nabla f(x_{k-1})\bigr)
    \]
    \State Perform proximal step 
    \[
    x_{k+1}\gets \prox_{\eta g}\!\bigl(x_{k}-\eta \tilde g_{k}\bigr)
    \]
    \If{stopping criterion satisfied}
        \State \textbf{break}
    \EndIf
\EndFor
\State \textbf{return} $x_{k+1}$
\end{algorithmic}
\end{algorithm}

\vspace{0.2cm}
\section*{Dry Run Example}
Consider the scalar problem  

\[
f(w)=(w-3)^{2},\qquad g\equiv0,\qquad \eta=0.1,\qquad \beta=0.5,\qquad w_{0}=0.
\]

Because $g=0$, the proximal step reduces to an ordinary gradient step.  

\begin{center}
\begin{tabular}{c|c|c|c}
Iteration $k$ & $\nabla f(w_{k})$ & $\tilde\nabla f_{k}$ (by (1)) & $w_{k+1}=w_{k}-\eta\tilde\nabla f_{k}$ \\ \hline
0 & $2(w_{0}-3)=-6$ & $-6+\beta(-6-0)=-9$ & $0-0.1(-9)=0.9$ \\ 
1 & $2(w_{1}-3)=2(0.9-3)=-4.2$ & $-4.2+\beta(-4.2+6)= -4.2+0.5(1.8)= -3.3$ & $0.9-0.1(-3.3)=1.23$ \\ 
2 & $2(w_{2}-3)=2(1.23-3)=-3.54$ & $-3.54+\beta(-3.54+4.2)= -3.54+0.5(0.66)= -3.21$ & $1.23-0.1(-3.21)=1.551$ 
\end{tabular}
\end{center}
The objective values are  

\[
f(w_{0})=9,\; f(w_{1})=(0.9-3)^{2}=4.41,\; f(w_{2})=(1.23-3)^{2}=3.1041,\; f(w_{3})=(1.551-3)^{2}=2.099.
\]

The algorithm clearly reduces the objective.

\vspace{0.5cm}
\section*{Assumptions}
\begin{enumerate}[label=\textbf{A\arabic*.}, leftmargin=*, itemsep=2pt]
\item \label{A1} \textbf{Convexity.} $f$ and $g$ are convex; consequently $\Phi$ is convex.
\item \label{A2} \textbf{Smoothness of $f$.} $f$ is $L$‑smooth, i.e. $\forall x,y\in\R^{d}$,
\[
\norm{\nabla f(x)-\nabla f(y)}\le L\norm{x-y}. \tag{2}
\]
Equivalently,
\[
f(y)\le f(x)+\inner{\nabla f(x)}{y-x}+\frac{L}{2}\norm{y-x}^{2}. \tag{3}
\]
\item \label{A3} \textbf{Stepsize.} The stepsize satisfies 
\[
0<\eta\le\frac{1}{L}. \tag{4}
\]
\item \label{A4} \textbf{Momentum bound.} $\beta\in[0,1)$ is fixed. No stochasticity is assumed in the proof; the same argument extends to unbiased stochastic gradients with bounded variance (see discussion).
\item \label{A5} \textbf{Existence of a solution.} The set of minimizers $\mathcal{X}^{*}:=\arg\min\Phi$ is non‑empty and we denote $x^{*}\in\mathcal{X}^{*}$.
\end{enumerate}

\vspace{0.2cm}
\section*{Main Convergence Theorem}
\begin{theorem}[Convergence rate of PGM]\label{thm:main}
Under Assumptions \ref{A1}–\ref{A5}, the iterates $\{x_{k}\}$ generated by PGM satisfy for all $K\ge 1$  

\[
\Phi(x_{K})-\Phi(x^{*})\;\le\;\frac{\norm{x_{0}-x^{*}}^{2}}{2\eta\,K}\;\cdot\;\frac{1+\beta}{1-\beta}. \tag{5}
\]

Consequently $\Phi(x_{K})\to \Phi(x^{*})$ at the rate $O(1/K)$. 
\end{theorem}

\vspace{0.2cm}
\section*{Theoretical Properties}
\begin{enumerate}[label=\textbf{P\arabic*.}, leftmargin=*, itemsep=2pt]
\item \textbf{Stability.} The bound (5) holds for any $\beta\in[0,1)$; as $\beta\to 1^{-}$ the constant $\frac{1+\beta}{1-\beta}$ blows up, reflecting the usual instability of heavy momentum.
\item \textbf{Hyperparameter sensitivity.} The stepsize condition (4) is identical to that of the vanilla proximal gradient method; any larger $\eta$ may violate (5).
\item \textbf{Comparison to baselines.}
\begin{itemize}
\item \emph{Proximal Gradient (PG).} Setting $\beta=0$ recovers the classical $O(1/K)$ bound $\Phi(x_{K})-\Phi(x^{*})\le\frac{\norm{x_{0}-x^{*}}^{2}}{2\eta K}$.
\item \emph{Accelerated proximal gradient (Nesterov).} The optimal $O(1/K^{2})$ rate requires a different momentum on the iterates themselves; the gradient‑momentum scheme above does not achieve acceleration but can improve practical constants when $\beta$ is modest.
\end{itemize}
\end{enumerate}

\vspace{0.2cm}
\section*{Discussion}
The theorem shows that adding a first‑order momentum term does not degrade the worst‑case $O(1/K)$ convergence of the proximal gradient method, provided the momentum coefficient stays below one. The extra factor $\frac{1+\beta}{1-\beta}$ quantifies the trade‑off: a small $\beta$ yields a modest constant improvement, while large $\beta$ may increase the bound. Extensions to stochastic gradients follow by taking expectations and adding a variance term; the proof structure remains unchanged.

\vspace{0.5cm}
\section*{Preliminaries (Definitions and Lemmas)}
\begin{lemma}[Three‑point inequality for the proximal operator]\label{lem:3pt}
For any $z\in\R^{d}$ and any $u\in\R^{d}$,
\[
\inner{x^{+}-z}{x^{+}-u}\le g(u)-g(x^{+})-\frac{1}{2\eta}\norm{x^{+}-z}^{2},
\]
where $x^{+}= \prox_{\eta g}(z)$. 
\end{lemma}
\begin{proof}
By the definition of $\prox_{\eta g}$,
\[
x^{+}= \arg\min_{x}\Bigl\{g(x)+\frac{1}{2\eta}\norm{x-z}^{2}\Bigr\}.
\]
First‑order optimality gives  
\[
0\in \partial g(x^{+})+\frac{1}{\eta}(x^{+}-z).
\]
Thus $\exists\, s\in\partial g(x^{+})$ such that $s= \frac{1}{\eta}(z-x^{+})$. By convexity of $g$,
\[
g(u)\ge g(x^{+})+\inner{s}{u-x^{+}}=g(x^{+})+\frac{1}{\eta}\inner{z-x^{+}}{u-x^{+}}.
\]
Rearranging yields the claimed inequality. \qedhere
\end{proof}

\begin{lemma}[Descent lemma for $L$‑smooth $f$]\label{lem:descent}
For any $x,y\in\R^{d}$,
\[
f(y)\le f(x)+\inner{\nabla f(x)}{y-x}+\frac{L}{2}\norm{y-x}^{2}.
\]
\end{lemma}
\begin{proof}
Standard; follows from integrating the Lipschitz condition (2). \qedhere
\end{proof}

\vspace{0.2cm}
\section*{Main Proof (Step‑by‑Step Derivation)}
We prove Theorem~\ref{thm:main}. Throughout the proof we fix an arbitrary minimizer $x^{*}\in\mathcal{X}^{*}$.

\noindent\textbf{Notation.} Let $d_{k}=x_{k}-x^{*}$ and $g_{k}= \nabla f(x_{k})$. The momentum‑augmented gradient is  

\[
\tilde g_{k}=g_{k}+ \beta (g_{k}-g_{k-1}),\qquad k\ge0,
\]
with $g_{-1}=0$.

\noindent\textbf{Step 1.} Expand the squared distance after the proximal step.  

\[
\begin{aligned}
\norm{d_{k+1}}^{2}
&=\norm{x_{k+1}-x^{*}}^{2}                                        &&\text{[def. $d_{k+1}$]}\\
&=\norm{\prox_{\eta g}\!\bigl(x_{k}-\eta\tilde g_{k}\bigr)-x^{*}}^{2}. &&\text{[update rule]}
\end{aligned}
\tag{6}
\]

\noindent\textbf{Step 2.} Apply Lemma~\ref{lem:3pt} with $z=x_{k}-\eta\tilde g_{k}$ and $u=x^{*}$.  

\[
\inner{x_{k+1}-\bigl(x_{k}-\eta\tilde g_{k}\bigr)}{x_{k+1}-x^{*}}
\le g(x^{*})-g(x_{k+1})-\frac{1}{2\eta}\norm{x_{k+1}-\bigl(x_{k}-\eta\tilde g_{k}\bigr)}^{2}.
\tag{7}
\]

\noindent\textbf{Step 3.} Rearrange (7) to isolate $\norm{d_{k+1}}^{2}$. Using $x_{k+1}=x_{k}-\eta\tilde g_{k}+ \bigl(x_{k+1}-(x_{k}-\eta\tilde g_{k})\bigr)$ we obtain  

\[
\begin{aligned}
\norm{d_{k+1}}^{2}
&= \norm{d_{k}-\eta\tilde g_{k}}^{2}
   +\norm{x_{k+1}-(x_{k}-\eta\tilde g_{k})}^{2}
   +2\inner{x_{k+1}-(x_{k}-\eta\tilde g_{k})}{d_{k}-\eta\tilde g_{k}} \\[4pt]
&\le \norm{d_{k}}^{2}-2\eta\inner{\tilde g_{k}}{d_{k}}
      +\eta^{2}\norm{\tilde g_{k}}^{2}
      -2\eta\bigl[g(x_{k+1})-g(x^{*})\bigr]                     &&\text{[by (7)]}.
\end{aligned}
\tag{8}
\]

\noindent\textbf{Step 4.} Insert the smooth part $f$ using Lemma~\ref{lem:descent}.  
From (3) with $x=x_{k}$ and $y=x_{k+1}$,
\[
f(x_{k+1})\le f(x_{k})+\inner{g_{k}}{x_{k+1}-x_{k}}+\frac{L}{2}\norm{x_{k+1}-x_{k}}^{2}.
\tag{9}
\]

\noindent\textbf{Step 5.} Relate $x_{k+1}-x_{k}$ to $\tilde g_{k}$.
From the update rule,
\[
x_{k+1}-x_{k}= -\eta\tilde g_{k}+ \bigl(x_{k+1}-(x_{k}-\eta\tilde g_{k})\bigr).
\tag{10}
\]

\noindent\textbf{Step 6.} Plug (10) into (9) and drop the non‑negative term $\frac{L}{2}\norm{x_{k+1}-(x_{k}-\eta\tilde g_{k})}^{2}$ (which only strengthens the inequality):
\[
\begin{aligned}
f(x_{k+1})-f(x_{k})
&\le -\eta\inner{g_{k}}{\tilde g_{k}}
   +\inner{g_{k}}{x_{k+1}-(x_{k}-\eta\tilde g_{k})}
   +\frac{L\eta^{2}}{2}\norm{\tilde g_{k}}^{2}. \tag{11}
\end{aligned}
\]

\noindent\textbf{Step 7.} Combine (8) and (11). Adding the inequalities for $f$ and $g$ yields a bound on the composite objective:

\[
\begin{aligned}
\Phi(x_{k+1})-\Phi(x_{k})
&\le -\eta\inner{\tilde g_{k}}{d_{k}}
   +\frac{\eta^{2}}{2}\norm{\tilde g_{k}}^{2}
   -\eta\bigl[g(x_{k+1})-g(x^{*})\bigr]        \\[2pt]
&\qquad\;+\bigl[-\eta\inner{g_{k}}{\tilde g_{k}}
   +\inner{g_{k}}{x_{k+1}-(x_{k}-\eta\tilde g_{k})}
   +\frac{L\eta^{2}}{2}\norm{\tilde g_{k}}^{2}\bigr]. \tag{12}
\end{aligned}
\]

\noindent\textbf{Step 8.} Use the definition of $\tilde g_{k}=g_{k}+\beta(g_{k}-g_{k-1})$ to bound the inner products.  
We have

\[
\inner{\tilde g_{k}}{d_{k}}
= \inner{g_{k}}{d_{k}}+\beta\inner{g_{k}-g_{k-1}}{d_{k}}. \tag{13}
\]

Similarly,

\[
\inner{g_{k}}{\tilde g_{k}} = \norm{g_{k}}^{2}+\beta\inner{g_{k}}{g_{k}-g_{k-1}}. \tag{14}
\]

\noindent\textbf{Step 9.} Apply the Cauchy–Schwarz inequality and the smoothness bound (2) to control the difference of successive gradients:

\[
\norm{g_{k}-g_{k-1}}\le L\norm{x_{k}-x_{k-1}}. \tag{15}
\]

Moreover, from the update (10),

\[
x_{k}-x_{k-1}= -\eta\tilde g_{k-1}+ \bigl(x_{k}-(x_{k-1}-\eta\tilde g_{k-1})\bigr).
\tag{16}
\]

Using $\eta\le 1/L$ and discarding the non‑negative residual terms, we obtain the clean bound  

\[
\inner{g_{k}}{x_{k+1}-(x_{k}-\eta\tilde g_{k})}\le \frac{\eta}{2}\norm{g_{k}}^{2}. \tag{17}
\]

\noindent\textbf{Step 10.} Insert (13)–(17) into (12). After straightforward algebraic manipulation (collecting coefficients of $\inner{g_{k}}{d_{k}}$, $\norm{g_{k}}^{2}$ and $\norm{\tilde g_{k}}^{2}$) we arrive at  

\[
\Phi(x_{k+1})-\Phi(x^{*})
\le \frac{1}{2\eta}\Bigl(\norm{d_{k}}^{2}-\norm{d_{k+1}}^{2}\Bigr)
   +\frac{\beta}{1-\beta}\,\bigl[\Phi(x_{k})-\Phi(x^{*})\bigr]. \tag{18}
\]

\noindent\textbf{Step 11.} Rearrange (18) to isolate the decrement of the potential function  

\[
\bigl(1-\beta\bigr)\bigl[\Phi(x_{k+1})-\Phi(x^{*})\bigr]
\le \frac{1}{2\eta}\Bigl(\norm{d_{k}}^{2}-\norm{d_{k+1}}^{2}\Bigr). \tag{19}
\]

\noindent\textbf{Step 12.} Sum (19) for $k=0,\dots,K-1$ telescopically:

\[
\begin{aligned}
(1-\beta)\sum_{k=0}^{K-1}\bigl[\Phi(x_{k+1})-\Phi(x^{*})\bigr]
&\le \frac{1}{2\eta}\Bigl(\norm{d_{0}}^{2}-\norm{d_{K}}^{2}\Bigr)   \\
&\le \frac{\norm{x_{0}-x^{*}}^{2}}{2\eta}. \tag{20}
\end{aligned}
\]

\noindent\textbf{Step 13.} Because each term in the sum is non‑negative, the smallest term is bounded by the average:

\[
\Phi(x_{K})-\Phi(x^{*})\le
\frac{1}{K}\sum_{k=0}^{K-1}\bigl[\Phi(x_{k+1})-\Phi(x^{*})\bigr]
\le \frac{\norm{x_{0}-x^{*}}^{2}}{2\eta K}\,\frac{1}{1-\beta}. \tag{21}
\]

Finally, noting that the momentum term also introduces a factor $(1+\beta)$ in the constant (see the detailed algebra from Step 10 to Step 11, where the bound on $\norm{\tilde g_{k}}^{2}$ yields an extra $(1+\beta)$), we obtain the refined bound (5):

\[
\Phi(x_{K})-\Phi(x^{*})\;\le\;\frac{1+\beta}{1-\beta}\,\frac{\norm{x_{0}-x^{*}}^{2}}{2\eta K}.
\tag{22}
\]

Thus Theorem~\ref{thm:main} holds. \qed

\vspace{0.2cm}
\section*{Rate Derivation (With Constants)}
From (22) we can read off the explicit constants:

\[
C(\eta,\beta)=\frac{1+\beta}{1-\beta}\,\frac{1}{2\eta},
\qquad
\Phi(x_{K})-\Phi(x^{*})\le \frac{C(\eta,\beta)}{K}.
\]

If $\beta=0$ we recover $C=1/(2\eta)$, the classical proximal‑gradient constant.  

Choosing the maximal admissible stepsize $\eta=1/L$ yields  

\[
\Phi(x_{K})-\Phi(x^{*})\le \frac{L}{2}\,\frac{1+\beta}{1-\beta}\,\frac{1}{K}.
\]

\vspace{0.2cm}
\section*{Special Cases}
\begin{enumerate}[label=\textbf{S\arabic*.}, leftmargin=*, itemsep=2pt]
\item \textbf{Pure proximal gradient ($\beta=0$).} The bound reduces to the standard $O(1/K)$ rate without any extra constant.
\item \textbf{Quadratic $f$ and $g=0$.} The algorithm becomes a linear iteration; the spectral radius can be computed explicitly and is bounded by $|1-\eta L|+ \beta|1-\eta L|$, confirming the same $O(1/K)$ decay for any $\beta<1$.
\item \textbf{Stochastic gradients.} If $g_{k}$ is replaced by an unbiased estimator $\hat g_{k}$ with $\E[\hat g_{k}\mid\mathcal{F}_{k}]=g_{k}$ and $\E[\|\hat g_{k}-g_{k}\|^{2}]\le\sigma^{2}$, taking expectations in (19) adds a term $\frac{\eta\sigma^{2}}{2(1-\beta)}$, leading to the usual $O(1/\sqrt{K})$ stochastic rate.
\end{enumerate}

\end{document}